% Capability hierarchy: colored name prefixes; neutral titles and captions.
% Data colors inside the plots retain their original meanings.
\usepackage{xcolor}
\definecolor{capSequencePriors}{HTML}{9C7524}
\definecolor{capFunctionalPrediction}{HTML}{2463A6}
\definecolor{capBiologicalContext}{HTML}{19836E}
\definecolor{capGeneticInteractions}{HTML}{7956A5}
\definecolor{capGenomicDesign}{HTML}{BE4D35}
\definecolor{capMutationEffectsAndContext}{HTML}{19836E}
\definecolor{capEditExecution}{HTML}{526674}
\definecolor{capPersonalPrediction}{HTML}{AE4275}

% Use descriptive category keys in every \cap... command.
% Example: \caphighlight{FunctionalPrediction}, \capfigure{GenomicDesign}.
\newcommand{\caplabelSequencePriors}{Sequence priors}
\newcommand{\caplabelFunctionalPrediction}{Functional prediction}
\newcommand{\caplabelBiologicalContext}{Biological context}
\newcommand{\caplabelGeneticInteractions}{Genetic interactions}
\newcommand{\caplabelGenomicDesign}{Genomic design}
\newcommand{\caplabelMutationEffectsAndContext}{Mutation Effects and Context}
\newcommand{\caplabelEditExecution}{Edit Execution}
\newcommand{\caplabelPersonalPrediction}{Personal Prediction}
% Compatibility aliases for older snippets and archived figure/table generators.
\colorlet{capDesignDecisions}{capGenomicDesign}
\let\caplabelDesignDecisions\caplabelGenomicDesign
\colorlet{capPrior}{capSequencePriors}
\let\caplabelPrior\caplabelSequencePriors
\colorlet{capSimulate}{capFunctionalPrediction}
\let\caplabelSimulate\caplabelFunctionalPrediction
\colorlet{capContext}{capBiologicalContext}
\let\caplabelContext\caplabelBiologicalContext
\colorlet{capCompose}{capGeneticInteractions}
\let\caplabelCompose\caplabelGeneticInteractions
\colorlet{capControl}{capGenomicDesign}
\let\caplabelControl\caplabelGenomicDesign
\colorlet{capResponse}{capMutationEffectsAndContext}
\let\caplabelResponse\caplabelMutationEffectsAndContext
\colorlet{capExecution}{capEditExecution}
\let\caplabelExecution\caplabelEditExecution
\colorlet{capPersonal}{capPersonalPrediction}
\let\caplabelPersonal\caplabelPersonalPrediction

\newcommand{\caplabel}[1]{\csname caplabel#1\endcsname}

% Section-name prefixes and vertical figure classification indices share a palette.
\DeclareRobustCommand{\capname}[1]{\textcolor{cap#1}{\caplabel{#1}}}
\DeclareRobustCommand{\capheading}[2]{\capname{#1}: #2}
\newcommand{\caplead}[2]{\textbf{#2}}
\pdfstringdefDisableCommands{%
  \def\capname#1{\caplabel{#1}}%
  \def\capheading#1#2{\caplabel{#1}: #2}}

% Capture the whole multi-panel figure, then place one full-height category index.
\newsavebox{\capPlotBox}
\newsavebox{\capLabelBox}
\newlength{\capPlotHeight}
\newlength{\capUpperRule}
\newlength{\capLowerRule}
\newcommand{\capfigure}[1]{%
  \begin{lrbox}{\capPlotBox}%
  \begin{minipage}[b]{\dimexpr\linewidth-16pt\relax}\centering}
\newcommand{\caprender}[1]{%
  \end{minipage}\end{lrbox}%
  \sbox{\capLabelBox}{\rotatebox{90}{%
    \textcolor{cap#1}{\fontfamily{phv}\fontseries{m}\fontshape{sl}%
    \fontsize{7}{8}\selectfont \caplabel{#1}}}}%
  \setlength{\capPlotHeight}{\dimexpr\ht\capPlotBox+\dp\capPlotBox\relax}%
% Fixed physical length, matching the accepted main Context figure.
  \setlength{\capUpperRule}{15.07bp}%
  \setlength{\capLowerRule}{\dimexpr\capPlotHeight-\capUpperRule-\ht\capLabelBox-\dp\capLabelBox-11pt\relax}%
  \par\noindent\hbox to\linewidth{\hfil%
    \raisebox{\dp\capPlotBox}{\usebox{\capPlotBox}}\hspace{6pt}%
    \vbox{\offinterlineskip%
      \hbox to8pt{\hfil\textcolor{cap#1!80!white}{\rule{.45pt}{\capUpperRule}}\hfil}%
      \vskip5.5pt%
      \hbox to8pt{\hfil\usebox{\capLabelBox}\hfil}%
      \vskip5.5pt%
      \hbox to8pt{\hfil\textcolor{cap#1!80!white}{\rule{.45pt}{\capLowerRule}}\hfil}}%
    \hfil}\par}
% Captions stay entirely black.
% A compact, consistent gap for classified figures; table captions are unchanged.
\newcommand{\capcaption}[2]{\caprender{#1}%
  \setlength{\abovecaptionskip}{4pt}\caption{#2}}

% Capability taxonomy: black lettering on a restrained category-colored highlight.
\DeclareRobustCommand{\caphighlight}[1]{%
  \begingroup\setlength{\fboxsep}{1.4pt}%
  \colorbox{cap#1!17!white}{\textcolor{black}{\caplabel{#1}:}}\endgroup}
\pdfstringdefDisableCommands{\def\caphighlight#1{\caplabel{#1}}}
